\chapter{Discussion}
\label{ch:discussion}

% For this thesis the research questions that were stated are:

% How would the participants describe the experience of using voice input in the presented
% game environment?

% My answer

% How do users feel about using the voice modality compared to traditional input?

% My answer 




% [Finding from Results] suggests that [interpretation].
% So what does that mean? So why does that matter?

\section{A - Prior experiences carry-over effect}
The observed behaviour pattern of players transferring their experience from the keyboard version into the voice experience suggests that participants do not approach the new input modalities in a vacuum. Instead, they use their prior experiences as reference points - be it keyboard navigation, voice assistants or mental models they have about human-computer interactions. 
This seems to be consistent with the found HCI literature on mental model transfer, where users continuously map existing knowledge and behaviour patterns into new input modalities \textcite{harneskExploringVoiceInteractions2019}.


This effect persisted even after the test order was reversed. The players who hadn't seen the keyboard version still used their past experiences to shape their interaction with the game. This implies that the modality transfer is not purely an effect of the A/B testing sequence.

This suggests that the way players felt about the voice input was already influenced by their prior experiences and biases. Therefore, the comparison that they were making between the voice-driven and keyboard-driven versions of the experience was not neutral. 





\section{B - Perception of control}

The differing reactions to open-ended voice input suggest that people view control differently. 
Findings suggest that for some people, control might be rooted in the ability to express their wishes openly. 
Meanwhile, for other players, being in control is about the outcome of their commands being what they want.
The open-ended nature of input made some people feel that they could ask for anything they wanted. That in itself seems to be the source of their feeling of control. 
On the other hand, the voice processing system refusals led some participants to feel less in control compared to using a keyboard.
Existing literature \textcite{cornelioSenseAgencyEmerging2022} supports this division into expression-based and outcome-based control perception.


% -WORK IN PROGRESS SECTION-

% finding: different emotional reaction 


% So what? Same feature - different results  

% People don't look at the same game in a same way. 
% Intrigue of mist or being scared by it

% Feelign more in crontrol cause I can ask for anything
% but I might have less of a tollerance for lacking responses. 

% Compare to RQ: Players felt visible emotions when playing the voice driven game, and there were different from the keyboard version


% ---

% The scope of how diverese the responses .. 


% The range of emotional statates thtat were produced by the players as a result of the open-ened nature of voice input suggests that not all participants percieve agency and feeling in control in the same way. 

% The non-deterministic nature of emotions that were produced during the voice-driven interaction suggests that pariticipants may react diffrently to the open-ended input and that their emotional tollerance to refusals vary. 

% The fact that some players reported a feeling of a higher agency compared to others despite both group encoutnering refusals from the game suffests that for some participants the ability to issue a command of their own choosing constifuted to that feeling of higher agency.

% For other participants the response gap - game not fufiling their commands was what made them feel like they're lacking agency and control over the experience.

% This suggests that some participants understand control as something rotted in in their intentions, while for the other group it's embodied by the results of their wishes. 





% Comparison of presented technologies/methods/projects

% Critical discussion / comparison of approaches

% Which methods are most commonly used, and why?

% Overview / summary of the literature reviewed, organised into a meaningful categorisation / characterisation


\section{C - The relationship between command language and the game experience}

% divergent forms -> people percieve AI differ
% relationship -> emotions

% people will form their own relationships with a voice input

% Why people speak differently? - results 

% Why does how you speak change how you feel? 
% Why is this important? 

% The divergent forms of relation that participants had with the game suggest that the voice-driven input had a social layer to it that was absent in the keyboard version.

% The divergent responses ... suggests that the voice driven input can change the emotions that particiapnts feel based 

% ... suggests that the human-computer relationship type has a direct impact on player emotions. 

% ... suggests that people form differing relationships with the voice input system based on their pre-existing notions about human computer interaction resulting in variation in emotional responses. 


% --- 
% The fact that people used differing interaction styles, some cooperative, yet another ones seeking permission, showcases that people may naturally drift towards a personal style of interaction if left uninstructed on how to communicate, be it issuing commands, asking for permission, speaking in first person or addressing the game like a cooperative partner.


The fact that people used differing interaction styles, showcases that people may naturally drift towards a personal style of interaction if left uninstructed on how to communicate. 

This suggests that participants are not treating the voice input as a deterministic input method; they view it as an input modality with an inherent non-linear social element. This aligns with the existing literature \cite{interactiondesignfoundationWhatAreVoice2025}, which states that voice communication plays a crucial role in human interaction. It goes further by claiming that people can't fully disregard the social expectations and frameworks of voice interaction, including the different ways of framing commands. 

The difference in emotional response in those interaction styles seems to come from how people naturally view power dynamics in social interactions. 
In this study, users who ask for permission feel less in control of their actions compared to those who simply express their commands. This is reflective of traditional game design in which participants report feeling more in charge when they  directly control the player character compared to less direct forms of player control \cite{kiiskiVoiceGamesHistory2020}.

% This discussed layer of emotions seems to be absent in the keyboard navigated experience further suggesting that the open-ended nature of voice input is responsible for those emotions in this case study.
% The fact that the above emotional layer was absent in the keyboard version shows that the voice driven input could make the game feel closer to a social interaction compared to keyboard input which doesn't force this social dynamic. 

% This emotional disparity between the two input method suggests that participants feel vastly different when playing both versions as a result of the forced social element of their interaction.

% relationship with game diffient
% power dynamics
% communication social layer

% keyboard is predictable
% voice carries one extra element that makes it social by definition 

Having this social interaction layer suggests that participants feel emotions that are a direct result of the input itself. It feels different to command a person to open the doors than top press a button to do it, since speaking carries that social element to that pressing a button doesn't. 


% Having this social interaction layer suggests that participants feel different emotions depending on their communication style, in turn making their 


% This social element in turn is what's responsible for the creation of the differing interaction patterns. 


% Participants seem to follow into two main categories of relationships with the voice input system. First one being asking for permission eg. "Can we go there?/Can I go there?" and the second one treating the system as a partner "Let's go there/I want to go there".

% This connects to exsit

\section{D - Emotional impact of technical limitations}

% why does this matter? People seem to respond differently 

% voice inroduces friction in a way that doesn't happen on keyboard 

% Higher mental load 

% Varied response - tolenrance 

% The range of emotional responses to the friction present in game suggests that participants might treat same pain points differently due to their personal tolerance to friction.

The fact that people often based their personal preference between the game versions based on the degree of frustration that they felt, combined with the varied responses that they had to similar pain points, suggests two things. Firstly, the tolerance for friction differs significantly between players, but frustration is also a crucial emotion surrounding game preference. 

% All of the participants played the same version of the game in each case encountering similar errors and bugs. The 
% The emotional response differential 


Given that some participants treated friction in-game as a source of major pain while others looked at it as a mere curiosity, it can be hypothesised that people might be affected differently by the same pain points. 

% Looking at the data about the participants, there doesn't seem to be any direct cause for this divide, as markers such as gaming experience or familiarity with voice assistants were not linked as predictors for this divide.

Looking at the data about the participants, the cause of this divide is unclear. The data points gathered during the testing, such as gaming experience or English level, are not reliable predictors. 
Existing literature indicates that both in the cases of voice-related frustration and the feelings of overload, the personal response varies between individuals, proving this finding but not offering an answer to why that is \textcite{harneskExploringVoiceInteractions2019}.



% Be it increased mental load, the errors in voice recognition or the bugs present in game players had varying responses.
This implies that voice-based input systems can leave can leave some participants feeling more irritated, frustrated, overloaded and angry compared to keyboard-driven input. 

When preferring the keyboard version, participants often made their choice based on comfort that they felt, not the fun and excitement they had experienced. This implies that players can value comfort above those emotions when playing games.




% This implies that when designing and testing games, any potential pain points noted by participants should not be dismissed purely because other players didn't feel in the same way.  







% The presence of frustration and cognitive load in some particiapnts both indicate that the voice-driven interactions increase the surface of potential pain points in games which is important from both the perspectives of the players but also the game developers planning to use voice input.




\section{E - Impact of using voice on player immersion}

% so what? Why is this important? Why does this matter?

% forced immersion wasn't wanted all the time
% Based on the responses the game felt more real 

% The split is interesting becuase 
% the definition says that immersion can be measured - we did, and we found out that it varies a lot. This is interesting. 

% Reasons for why something is immersive differ a lot - also super interesting. Personal immersion triggers. 

% People not preffering immerison. Dommination of friction above all else. 

% Friction should stop immersion, turns out sometimes it doesn't !! for some people it does 

% My findings complicte the definition implying that it's one of the most crucial metrics, it is but there's a deeper more fundamental metric here.
% Friction as a more important metric in game assessment, beats immersion. Players chose their prefference not based on immersion but comfort. 
% That's why people didn't like the immersion, it came at the cost of comfort.

% The fact that the voie inpput modality had a bipolar impact on the players immersion suggests that each person has their own personal immersion triggers that often don't correlate between different people. 

% This comes at at surprise considering that 

% Throughout the study each person gave their own reasoning behind why they found voice input modality more or less immersive. 

% In literature measuring immersion can be done both qualitatively and quantitavely. \todo{cite this}

% It should be noted that there are some general patters suggesting that for instance player frustrtion acts as an immersion suppressant which is universally agreed upon in the literature. \todo{cite this}
% Despite those universal markers the fact that participants offered very personal reasoning behind their immersion state showcases that it's difficult to create an experience that's universally immersive for all participants due to the varying reactions from the same stimuli. 


% The bipolar effect that voice input had on player immersion 


% Immersion and comfort at the same scale - that turned out to be false.

% During the study players often voiced their prefference for the keyboard version despite claiming that the voice-driven game is more immersive. 

% When aksed for their prefference particiapnts often 

% Going into this study the assumption was that friction would naturally break the feeling of immersion in participants. 
The fact that participants could claim both a higher degree of immersion in the voice-driven game and not choose that version, quoting frustration caused by friction, suggests that friction might not be a reliable predictor of immersion. 
The data suggests that players can feel deeply immersed even under high discomfort when playing the game. This was a surprising finding, even considering that literature supports a similar claim that players can reject a more immersive input system solely out of the lower efficiency of it \cite{allisonDesignPatternsVoice2018}. 


Another suggestion that can be made is the fact that players might all have their own personal factors that can make or break immersion, and that those variables are not universal across the sample size. 
This demonstrates that creating experiences that are immersive might prove to be challenging since immersion has been shown to be a complex and personal feeling. 
This correlates with the literature, which notes that the process of acquiring immersion while using voice interaction is deeply personal and that it varies between people, which goes hand in hand with the above findings \cite{harneskExploringVoiceInteractions2019}.



% It also shows that immersion is not a feeling that can be predicted based on the physical state, like pain as a result of trauma. It seems to be closer to a pers




% There is no cake that everyone can enjoy, but people will universally react to pain. 

Despite the mentioned complexity and friction, all the participants agreed on the fact that the input modality had a profound impact on their feeling of immersion in the game, be it positive or negative; it left a mark.